\chapter{Conclusion Générale et Perspectives}

\section{Conclusion Technique}
Le projet HR Analytics a permis de démontrer la puissance de l'écosystème Power BI dans la résolution de problématiques RH complexes. En migrant d'un système de reporting manuel vers une architecture décisionnelle structurée en Schéma en Étoile, l'organisation dispose désormais d'une source unique de vérité. La séparation modulaire entre l'ETL, la modélisation et la restitution garantit une robustesse technique et une facilité de maintenance sur le long terme.

\section{Bilan Collaboratif}
Sur le plan humain, cette réalisation est le fruit d'une synergie efficace entre quatre experts. L'utilisation d'un workflow Git basé sur des branches de fonctionnalités (\textit{feature branches}) a permis de paralléliser les tâches tout en assurant une intégration fluide des composants. Chaque membre a pu apporter sa pierre à l'édifice, de l'ingénierie des données brute jusqu'à la narration visuelle finale.

\section{Perspectives d'Évolution}
Bien que la solution actuelle réponde aux besoins descriptifs immédiats, plusieurs axes d'amélioration peuvent être envisagés pour passer d'une analyse descriptive à une analyse prescriptive :

\begin{itemize}
    \item \textbf{Analyse Prédictive (Machine Learning) :} Intégration de modèles de classification (ex: Random Forest ou Régression Logistique) pour calculer un \textit{Flight Risk Score} individuel, permettant d'identifier les départs probables avant qu'ils ne surviennent.
    \item \textbf{Sécurité à la ligne (RLS) :} Mise en place de la \textit{Row-Level Security} pour restreindre l'accès aux données sensibles. Par exemple, un chef de département ne pourrait voir que les données de ses propres collaborateurs.
    \item \textbf{Intégration de données externes :} Corrélation des données internes avec des benchmarks de salaires du marché pour affiner l'analyse de compétitivité.
    \item \textbf{Application Mobile :} Déploiement d'une version optimisée pour Power BI Mobile, permettant aux décideurs de consulter les KPIs clés en mobilité.
\end{itemize}

En conclusion, ce projet ne constitue pas seulement un outil de reporting, mais un véritable levier de transformation stratégique, plaçant la donnée au cœur de la gestion des talents de demain.
